% GENERATED FILE. Edit canonical lessons, then run npm run generate:books.
% canonical-source-hash: fnv1a64:b02f8440fbfdd5c5
% canonical-lessons: PA-R160-fall-again, PA-R160-teach-again, PA-R160-lift-again, PA-R160-jump-again

\chapter{Return to Four Familiar Action Words}
\label{ch:pa-four-familiar-action-words}
\hlchaptermodality{160}

\begin{chapteropening}
\textbf{By the end of this chapter:} \emph{I can retrieve the already-taught Punjabi verbs for falling, teaching, lifting, and jumping after a long gap, one at a time, without claiming a scored A1 mock.}

Recall the four familiar Chapter 148 actions from short cues without introducing new Punjabi or claiming scored writing credit.
\end{chapteropening}

\section[\texorpdfstring{\pa{ਡਿੱਗਣਾ} (ḍiggṇā)}{ਡਿੱਗਣਾ (ḍiggṇā)}]{An old word for falling}

\label{lesson:PA-R160-fall-again}



\begin{quote}
Hear \textbf{to fall}. {[}PAUSE 4s{]} Say the Punjabi word before the model: \textbf{\pa{ਡਿੱਗਣਾ}} (\emph{ḍiggṇā}).
\end{quote}

\subsection*{Your turn}
Look at \textbf{\pa{ਡਿੱਗਣਾ}} and say its familiar meaning. {[}PAUSE 4s{]} The meaning is \textbf{to fall}. Now hear \textbf{to fall} again, look away, and say \textbf{\pa{ਡਿੱਗਣਾ}} before the model. No sentence is needed yet.

\begin{tcolorbox}[breakable,colback=teal!4,colframe=teal!35!black,title={Before you move on}]
Hear \textbf{to fall} once more. {[}PAUSE 4s{]} Recall \textbf{\pa{ਡਿੱਗਣਾ}}. If it did not come back, listen to the model and retry once. This untimed word return is not part of a scored A1 mock.
\end{tcolorbox}
\section[\texorpdfstring{\pa{ਸਿਖਾਉਣਾ} (sikhāuṇā)}{ਸਿਖਾਉਣਾ (sikhāuṇā)}]{An old word for teaching}

\label{lesson:PA-R160-teach-again}



\begin{quote}
Hear \textbf{to teach}. {[}PAUSE 4s{]} Say the earlier Punjabi word before the model: \textbf{\pa{ਸਿਖਾਉਣਾ}} (\emph{sikhāuṇā}).
\end{quote}

\subsection*{Your turn}
Hear the model \textbf{\pa{ਸਿਖਾਉਣਾ}}. {[}PAUSE 4s{]} Recall its old meaning, \textbf{to teach}. Then hear \textbf{to teach} and say \textbf{\pa{ਸਿਖਾਉਣਾ}} once more before the spoken model. There is no new grammar here.

\begin{tcolorbox}[breakable,colback=teal!4,colframe=teal!35!black,title={Before you move on}]
One last cue: \textbf{to teach}. {[}PAUSE 4s{]} Recall \textbf{\pa{ਸਿਖਾਉਣਾ}} before the model. Retry a missed old word once; do not count this as a timed or scored writing attempt.
\end{tcolorbox}
\section[\texorpdfstring{\pa{ਚੁੱਕਣਾ} (cukkṇā)}{ਚੁੱਕਣਾ (cukkṇā)}]{An old word for lifting}

\label{lesson:PA-R160-lift-again}



\begin{quote}
Hear \textbf{to lift}. {[}PAUSE 4s{]} Recall \textbf{\pa{ਚੁੱਕਣਾ}} (\emph{cukkṇā}).
\end{quote}

\subsection*{Your turn}
Look at \textbf{\pa{ਚੁੱਕਣਾ}} and remember \textbf{to lift}. {[}PAUSE 4s{]} Look away. Hear \textbf{to lift} and say \textbf{\pa{ਚੁੱਕਣਾ}} before the spoken model. Stay with the old word; no new action is being taught.

\begin{tcolorbox}[breakable,colback=teal!4,colframe=teal!35!black,title={Before you move on}]
Hear \textbf{to lift} one final time. {[}PAUSE 4s{]} Say \textbf{\pa{ਚੁੱਕਣਾ}}. Listen and retry once if needed. This is an untimed vocabulary return, separate from the scored A1 writing tasks.
\end{tcolorbox}
\section[\texorpdfstring{\pa{ਟੱਪਣਾ} (ṭappṇā)}{ਟੱਪਣਾ (ṭappṇā)}]{An old word for jumping}

\label{lesson:PA-R160-jump-again}



\begin{quote}
Hear \textbf{to jump}. {[}PAUSE 4s{]} Say the old Punjabi word before the model: \textbf{\pa{ਟੱਪਣਾ}} (\emph{ṭappṇā}).
\end{quote}

\subsection*{Your turn}
Pause after each familiar meaning, then listen to the old word:

\begin{itemize}
\raggedright
  \item \textbf{to lift} — {[}PAUSE 4s{]} \textbf{\pa{ਚੁੱਕਣਾ}} (\emph{cukkṇā}).
  \item \textbf{to fall} — {[}PAUSE 4s{]} \textbf{\pa{ਡਿੱਗਣਾ}} (\emph{ḍiggṇā}).
  \item \textbf{to jump} — {[}PAUSE 4s{]} \textbf{\pa{ਟੱਪਣਾ}} (\emph{ṭappṇā}).
  \item \textbf{to teach} — {[}PAUSE 4s{]} \textbf{\pa{ਸਿਖਾਉਣਾ}} (\emph{sikhāuṇā}).
\end{itemize}

The shuffled order is the only new part of this practice.

\begin{tcolorbox}[breakable,colback=teal!4,colframe=teal!35!black,title={Before you move on}]
Hear \textbf{to teach}, \textbf{to jump}, \textbf{to fall}, then \textbf{to lift}. Pause after each cue and recall its Punjabi word before the model: \textbf{\pa{ਸਿਖਾਉਣਾ}}, \textbf{\pa{ਟੱਪਣਾ}}, \textbf{\pa{ਡਿੱਗਣਾ}}, \textbf{\pa{ਚੁੱਕਣਾ}}. Revisit only the old word you missed.
\end{tcolorbox}
